% ============================================================
%  LITERATURE REVIEW — Task Offloading in IoT–Fog–Cloud Systems
%  Deadline: April 15
%
%  HOW THIS FILE IS ORGANIZED FOR YOUR TEAM:
%  - YOUR 10 PAPERS  : Sections 1–5 (Introduction, Background,
%                      Evolution, State of the Art, Gaps)
%  - TEAMMATE SLOTS  : Look for  %% TODO-TEAMMATE-X blocks
%                      Each teammate adds their sections there
%  - READY TO SUBMIT : If teammates don't edit, it compiles
%                      cleanly as a complete standalone paper
% ============================================================

\documentclass[12pt,a4paper]{article}

% ── Packages ─────────────────────────────────────────────────────────────────
\usepackage[a4paper, margin=2.5cm, headheight=15pt]{geometry}
\usepackage{amsmath,amssymb}
\usepackage{graphicx}
\usepackage{booktabs}
\usepackage{array}
\usepackage{tabularx}
\usepackage{hyperref}
\usepackage[numbers,sort&compress]{natbib}
\usepackage{microtype}
\usepackage{parskip}
\usepackage{titlesec}
\usepackage{fancyhdr}
\usepackage{enumitem}
\usepackage{caption}
\usepackage{float}
\usepackage{tikz}
\usepackage{pgfplots}
\pgfplotsset{compat=1.18}
\usetikzlibrary{arrows.meta, positioning, shapes.geometric,
                fit, backgrounds, calc}

% ── Colours ──────────────────────────────────────────────────────────────────

% ── Hyperref ─────────────────────────────────────────────────────────────────
\hypersetup{
  colorlinks=false,
  hidelinks,
  pdftitle={Task Offloading in IoT-Fog-Cloud: A Literature Review},
}

% ── Headers ──────────────────────────────────────────────────────────────────
\pagestyle{fancy}
\fancyhf{}
\rhead{\small Task Offloading in IoT--Fog--Cloud: A Literature Review}
\lhead{\small \thepage}
\renewcommand{\headrulewidth}{0.4pt}

% ── Section formatting ────────────────────────────────────────────────────────
\titleformat{\section}
  {\large\bfseries}{\thesection}{1em}{}
\titleformat{\subsection}
  {\normalsize\bfseries}{\thesubsection}{1em}{}
\titleformat{\subsubsection}
  {\normalsize\bfseries}{\thesubsubsection}{1em}{}

% ── Custom environments ───────────────────────────────────────────────────────




% ── TODO box for team collaboration ──────────────────────────────────────────


% ── Column types ─────────────────────────────────────────────────────────────
\newcolumntype{L}[1]{>{\raggedright\arraybackslash}p{#1}}
\newcolumntype{C}[1]{>{\centering\arraybackslash}p{#1}}

% ── Helper command for team notes ─────────────────────────────────────────────
\newcommand{\teamtodo}[1]{%
  \begin{quote}\small\textit{\textbf{TEAM NOTE:} #1}\end{quote}
%
}

% =============================================================================
\begin{document}
% =============================================================================

% ── TITLE PAGE ────────────────────────────────────────────────────────────────
\begin{titlepage}
  \centering
  \vspace*{1.5cm}
  {\small\textcolor{midblue}{\textsc{Literature Review --- Distributed Systems Lab}}}\\[0.8cm]
  {\LARGE\bfseries\color{darkblue}
    Optimized Task Offloading in\\[0.3em]
    IoT--Fog--Cloud Systems:\\[0.4em]
    A Comprehensive Literature Review of\\[0.3em]
    Metaheuristic Optimization, Deep\\[0.3em]
    Reinforcement Learning, SDN, and\\[0.3em]
    Mobility-Aware Computation Offloading
  }\\[1.2cm]
  \rule{0.55\linewidth}{0.4pt}\\[1.0cm]
  {\normalsize
    \textbf{Keywords:} Task offloading, IoT, Fog computing, NSGA-II,\\
    Deep reinforcement learning, SDN, DAG, Vehicular networks,\\
    Mobility-aware offloading, Proactive handoff
  }\\[1.5cm]
  {\normalsize
    \textbf{Team:} imed eddine benmadi , 
  }\\[1.5cm]
  \rule{0.55\linewidth}{0.4pt}\\[0.8cm]
  {\normalsize\textcolor{midblue}{April 2026}}
  \vfill
  {\small\textit{Prepared for: Distributed Systems Lab --- Literature Review Assignment}}
\end{titlepage}
% ── ABSTRACT ─────────────────────────────────────────────────────────────────
\begin{abstract}
\noindent
The rapid proliferation of Internet of Things (IoT) devices in smart-city and
vehicular environments has created an urgent demand for low-latency, energy-efficient
computation at the network edge. Task offloading---the process of delegating
computation from resource-constrained IoT devices to fog or cloud servers---has
emerged as the central mechanism for satisfying these demands within the
three-tier IoT--Fog--Cloud hierarchy. This paper presents a structured literature
review of ten representative works spanning 2021--2025, covering five distinct
research directions: (i)~multi-objective metaheuristic optimization via NSGA-II
with Minimax Differential Evolution, (ii)~task pre-filtering through the
TOF-Broker Execution Cost framework, (iii)~software-defined networking for
programmable traffic management, (iv)~microservice-oriented DAG task modeling,
and (v)~mobility-aware offloading with deep reinforcement learning and proactive
handoff prediction. We identify the state of the art in each sub-domain, synthesize
the research trajectory, and formally define five open research gaps that no
published system has jointly resolved. Our analysis argues that the convergence of
offline Pareto optimization and online reinforcement learning---applied to DAG
workloads with GPS-driven mobility prediction---represents the most promising
direction for next-generation vehicular edge computing systems.

%% TODO-TEAMMATE-1: If your team wants to extend the abstract to include
%% a brief mention of your teammates' research directions (e.g., security,
%% federated learning, or whatever areas their 10 papers cover), add 2-3
%% sentences here. Otherwise leave as-is — it reads cleanly without it.
\end{abstract}

\vspace{0.4cm}\hrule\vspace{0.6cm}
\tableofcontents
\newpage

% =============================================================================
\section{Introduction}
\label{sec:intro}
% =============================================================================

The Internet of Things ecosystem is undergoing unprecedented growth, with tens of
billions of connected devices projected by 2030~\cite{jafari2021}. A growing
proportion of these devices operate in latency-sensitive, mobility-intensive
environments: smart vehicles performing real-time collision avoidance, autonomous
drones executing path-planning algorithms, and industrial robots running precision
control loops. These workloads share a critical characteristic---they generate
computation-intensive tasks that exceed the processing capacity of the originating
device yet must produce results within strict timing constraints measured in tens
of milliseconds.

Conventional cloud computing provides near-unlimited computational resources but
introduces WAN propagation delays of 50--200\,ms, rendering it unsuitable for
safety-critical vehicular applications~\cite{jafari2021}. Fog computing addresses
this by positioning intermediate compute nodes one or two network hops from IoT
devices, reducing latency to 1--10\,ms. The resulting \emph{three-tier hierarchy}
(IoT $\rightarrow$ Fog $\rightarrow$ Cloud) creates a fundamental optimization
problem: for each incoming task, which tier executes it, under what scheduling
constraints, and at what energy cost?

\emph{Task offloading}---the decision process routing computation across tiers---is
formally NP-hard~\cite{mashal2024drl}. At IoT scale, the search space is
intractable for exact algorithms. Five interconnected challenges define the current
research landscape:

\begin{enumerate}[leftmargin=*, label=\textbf{C\arabic*.}]
  \item \textbf{Boulder Problem.} Compute-heavy tasks routed naively to fog nodes
    consume disproportionate resources and cause queue cascade failures~\cite{nouri2024tof}.
  \item \textbf{Cold-Start Problem.} Deep reinforcement learning agents produce
    near-random policies during early training, causing service degradation before
    they converge~\cite{vfc_drl2024}.
  \item \textbf{Static Optimization Problem.} Evolutionary optimizers like NSGA-II
    converge over hundreds of generations---orders of magnitude too slow for
    millisecond real-time decisions~\cite{zhu2025dag}.
  \item \textbf{Network Congestion Problem.} Reactive SDN routing cannot respond
    fast enough to burst traffic events, causing packet loss and queuing
    delays~\cite{csfd2025,gasmi2024sdn}.
  \item \textbf{Mobility Problem.} In vehicular networks, devices traverse fog
    coverage zones faster than tasks complete, causing mid-computation handoff
    failures~\cite{mavto2024,paper5_mobility}.
\end{enumerate}

No single published system resolves all five challenges in a unified framework.
This literature review maps the research evolution from single-objective heuristics
to adaptive AI-driven systems, identifies the specific gaps that remain open,
and positions a proposed architecture---the Predictive Cloud-Native Mobile Edge
(PCNME)---as the natural convergence point of this trajectory.

\textbf{Paper organization.}
Section~\ref{sec:background} establishes technical background and formal
definitions. Section~\ref{sec:evolution} traces four generations of research
evolution. Section~\ref{sec:sota} surveys the state of the art in each
sub-domain with our ten primary references. Section~\ref{sec:gaps} formalizes
five open research gaps. Section~\ref{sec:proposed} describes the proposed
unified architecture. Section~\ref{sec:discussion} provides comparative
analysis. Section~\ref{sec:conclusion} concludes.

%% TODO-TEAMMATE-1 and TEAMMATE-2: You may add a sentence here referencing
%% the additional research directions covered in your respective sections
%% (Sections 8 and 9). Example: "Section 8 examines [your topic].
%% Section 9 addresses [other teammate's topic]."

% =============================================================================
\section{Background and Formal Definitions}
\label{sec:background}
% =============================================================================

\subsection{Three-Tier IoT--Fog--Cloud Architecture}

The three-tier model partitions computation infrastructure into:
\textbf{Layer~1} (IoT devices)---resource-constrained task generators with limited
CPU, memory, and battery;
\textbf{Layer~2} (fog nodes)---intermediate edge servers co-located with network
access points, offering low-latency compute within a coverage radius of 150--300\,m
in urban 5G deployments; and
\textbf{Layer~3} (cloud)---centralized high-throughput servers providing batch
computation at the cost of 20--50\,ms additional WAN propagation delay.

\subsection{Task Offloading Formulation}

Let $\mathcal{T} = \{t_1, t_2, \ldots, t_N\}$ be a batch of $N$ tasks arriving
from IoT devices. Each task $t_i$ is characterized by:
\begin{equation}
  t_i = \langle \ell_i,\; d_i,\; D_i,\; \mathcal{E}_i \rangle
\end{equation}
where $\ell_i$ is the computational demand in Million Instructions~(MI), $d_i$
is the input data size in kilobytes, $D_i$ is the end-to-end deadline in
milliseconds, and $\mathcal{E}_i$ is the device transmission energy in joules.

The offloading decision assigns each task to a destination
$x_i \in \{\text{FOG}_1, \ldots, \text{FOG}_K, \text{CLOUD}\}$.
The two primary optimization objectives are:
\begin{align}
  f_1(\mathbf{x}) &= \sum_{i=1}^{N} E_i(x_i)
    \quad \text{(total energy, minimize)} \label{eq:energy}\\
  f_2(\mathbf{x}) &= \sum_{i=1}^{N} L_i(x_i)
    \quad \text{(total latency, minimize)} \label{eq:latency}
\end{align}
subject to:
\begin{align}
  L_i(x_i) &\leq D_i \quad \forall\, i
    \quad \text{(deadline constraint)} \\
  \sum_{i:\,x_i = k} \ell_i &\leq C_k \quad \forall\, k
    \quad \text{(fog capacity constraint)}
\end{align}
where $C_k$ is the MIPS capacity of fog node $k$.

\subsection{Execution Cost and Boulder/Pebble Classification}

The \emph{Execution Cost}~(EC) metric~\cite{nouri2024tof} determines fog
feasibility:
\begin{quote}
\textbf{Definition (Execution Cost):}
\begin{equation}
  \mathrm{EC}(t_i) = \frac{\ell_i}{\text{fog\_MIPS}} \quad\text{(seconds)}
  \label{eq:ec}
\end{equation}
Task $t_i$ is a \emph{Pebble} if $\mathrm{EC}(t_i) < \theta$, and a
\emph{Boulder} if $\mathrm{EC}(t_i) \geq \theta$, where $\theta = 1.0$\,s
is the default threshold~\cite{nouri2024tof}.
\end{quote}


\subsection{Pareto Dominance and Non-Dominated Sorting}

\textbf{Definition (Pareto dominance):}
Solution $\mathbf{x}^{(a)}$ dominates $\mathbf{x}^{(b)}$, written
$\mathbf{x}^{(a)} \prec \mathbf{x}^{(b)}$, iff:
\begin{equation}
  f_k\!\left(\mathbf{x}^{(a)}\right) \leq f_k\!\left(\mathbf{x}^{(b)}\right)\;
  \forall k,\quad
  \exists j:\; f_j\!\left(\mathbf{x}^{(a)}\right) < f_j\!\left(\mathbf{x}^{(b)}\right)
\end{equation}

The Pareto front $\mathcal{P}^*$ is the set of all non-dominated solutions.
NSGA-II~\cite{jafari2021} approximates $\mathcal{P}^*$ by evolving a population
of candidate routing plans across multiple generations via selection, crossover,
and mutation operators.

\subsection{DAG Task Model}

Modern IoT applications decompose into chains of dependent processing
steps~\cite{li2025microservice}. The Directed Acyclic Graph (DAG) model
represents an application as $G = (\mathcal{V}, \mathcal{E})$ where node
$v_i \in \mathcal{V}$ is a microservice step with properties
$\langle \ell_i, d_i^{\text{in}}, d_i^{\text{out}} \rangle$,
and edge $(v_i, v_j) \in \mathcal{E}$ encodes data dependency: step $v_j$
cannot start until $v_i$ completes and transmits its output. Total application
latency over the critical path is:
\begin{equation}
  L_{\text{app}} = \sum_{v_i \in \text{critical path}}
  \!\!\Bigl[ T_{\text{exec}}(v_i) + T_{\text{transfer}}(v_i, v_{i+1}) \Bigr]
  \label{eq:dag_latency}
\end{equation}
where $T_{\text{transfer}} = 0$ when consecutive steps are co-located.

% =============================================================================
\section{Research Evolution: Four Generations}
\label{sec:evolution}
% =============================================================================

\subsection{Generation 1: Single-Objective Heuristics (pre-2021)}

Early task offloading research modeled the decision as binary: execute locally
or offload entirely to the cloud. Solutions used greedy algorithms, game theory,
and single-objective metaheuristics such as Particle Swarm Optimization and Ant
Colony Optimization~\cite{jafari2021}. These approaches optimized one metric while
ignoring the other, and operated on static topologies with no mobility modeling.

\textbf{Key limitation:} Single-objective formulations systematically underperform
because energy and latency conflict. Minimizing latency alone drives all tasks to
cloud (high energy); minimizing energy alone keeps tasks local (high latency). A
balanced solution requires explicit multi-objective treatment.

\subsection{Generation 2: Multi-Objective Evolutionary Algorithms (2021--2023)}

Jafari and Rezvani~\cite{jafari2021} established the multi-objective framing by
applying NSGA-II with Minimax Differential Evolution (MMDE) mutation to the
IoT--Fog--Cloud offloading problem. Their binary chromosome encodes routing
decisions for a task batch; solutions are evaluated on energy~\eqref{eq:energy}
and latency~\eqref{eq:latency} objectives simultaneously, producing Pareto fronts
that expose the full energy--latency trade-off space.

The MMDE operator~\cite{jafari2021} replaces standard random mutation with
directional differential evolution:
\begin{equation}
  \mathbf{V} = \mathbf{r}_1 + F \cdot (\mathbf{r}_2 - \mathbf{r}_3),
  \quad F = 0.5,\; \mathrm{CR} = 0.9
  \label{eq:mmde}
\end{equation}
where $\mathbf{r}_1, \mathbf{r}_2, \mathbf{r}_3$ are three distinct population
members. This improves convergence speed and solution robustness against
variable network conditions.

\textbf{Key limitation:} NSGA-II evaluates all tasks---including
compute-infeasible boulders---on fog nodes, wasting CPU cycles and causing queue
crashes. Additionally, hundreds-of-generations batch convergence is incompatible
with millisecond real-time requirements.

\subsection{Generation 3: Pre-Filtering and Hybrid Systems (2024)}

Nouri et al.~\cite{nouri2024tof} introduced the TOF-Broker atop NSGA-II.
By computing EC~\eqref{eq:ec} for each task before optimization, the broker
eliminates boulders from the search space, reducing problem size by 15--30\% and
achieving 12.18\% energy improvement over plain NSGA-II with negligible latency
overhead. Concurrently, the microservice research community~\cite{li2025microservice}
established that DAG-structured task models enable partial offloading---running
lightweight steps on fog while heavy steps execute on cloud---impossible under the
monolithic task model.

\textbf{Key limitation:} The resulting systems remain static. Pre-filtered
NSGA-II produces optimized batch routing plans that cannot adapt to changing
fog loads, varying network conditions, or moving devices.

\subsection{Generation 4: Adaptive AI-Driven Systems (2024--2025)}

The current generation replaces or augments offline optimization with online
learning agents. For vehicular settings, DRL-based offloading~\cite{vfc_drl2024}
demonstrated that DQN agents trained to minimize queue congestion outperform
greedy policies under high load. CSFD~\cite{csfd2025} embedded DRL inside the
SDN controller for proactive congestion-aware routing. MAVTO~\cite{mavto2024}
introduced GPS-based trajectory prediction to pre-position tasks at the
vehicle's destination fog zone before handoff. Zhu et al.~\cite{zhu2025dag}
combined DAG modeling, NSGA-II offline pre-training, and online meta-RL adaptation
in a single pipeline.

\textbf{Key limitation:} No system in this generation combines multi-objective
offline optimization, DRL online adaptation, SDN network management, DAG task
modeling, and proactive GPS-driven mobility---all five capabilities simultaneously.

% =============================================================================
\section{State of the Art}
\label{sec:sota}
% =============================================================================

\subsection{Multi-Objective Metaheuristic Optimization}

NSGA-II remains the dominant metaheuristic for multi-objective task offloading.
Jafari and Rezvani~\cite{jafari2021} established the canonical formulation: a
population of routing chromosomes (arrays of fog/cloud decisions per task) is
evolved using non-dominated sorting, crowding distance, SBX crossover, and MMDE
mutation~\eqref{eq:mmde}. Their evaluation on a 3-tier IoT--Fog--Cloud testbed
demonstrated that Pareto front solutions provide meaningful energy--latency
trade-off diversity, enabling operators to select operating points according to
their QoS budget.

The key innovation of MMDE over standard polynomial mutation is directionality:
the mutation vector $(r_2 - r_3)$ encodes population structure, steering mutations
toward productive search regions rather than sampling uniformly at random. This
produces faster convergence and more robust solutions under variable fog loads.

\subsection{Task Pre-Filtering: The TOF-Broker}

Nouri et al.~\cite{nouri2024tof} identified the boulder problem as the dominant
failure mode of Generation-2 systems and proposed the TOF-Broker to address it.
The broker applies EC classification~\eqref{eq:ec} to each incoming task before
any optimizer runs: boulders ($\mathrm{EC} \geq 1.0$\,s) are immediately routed
to cloud; pebbles ($\mathrm{EC} < 1.0$\,s) enter the NSGA-II optimization
pipeline. This achieves a 12.18\% energy reduction with only 0.38\% latency
overhead, while eliminating fog node crashes under heavy-task workloads.

However, the TOF-Broker introduces the \emph{Pebble Avalanche} problem: when
task arrival rates exceed $\sim$1,000 tasks/second, the pebble queue grows faster
than NSGA-II can drain it, eventually causing the same cascade failures the broker
was designed to prevent. This remains an open problem addressed in
Section~\ref{sec:gaps}.

\subsection{Software-Defined Networking for Fog Systems}

Gasmi and Ezzedine~\cite{gasmi2024sdn} conducted a systematic survey of SDN-based
fog offloading, establishing the theoretical foundation for programmable network
management in this domain. Their key finding: SDN's control/data plane
separation---placing routing intelligence in a centralized controller while
forwarding switches execute OpenFlow rules blindly---enables instant, globally
consistent traffic reprogramming from a single API call, unavailable in traditional
distributed routing.

This capability is critical for task offloading: when a large batch of pebbles
needs to be sent to the cloud (a Super-Task), the SDN controller can open a
dedicated high-priority lane on all intervening switches in under 5\,ms, bypassing
background telemetry traffic. However, Gasmi and Ezzedine~\cite{gasmi2024sdn}
also identify the centralized controller bottleneck: under high IoT device density,
the controller becomes overloaded with new-flow requests, adding 8--15\,ms latency
per flow---negating the benefit of SDN for latency-sensitive applications.

\subsection{Microservice Architecture and DAG Task Modeling}

Li et al.~\cite{li2025microservice} surveyed 136 studies (2020--2024) and produced
the first unified taxonomy bridging microservices orchestration and energy
efficiency across five resource management pillars: service placement, resource
provisioning, task scheduling and offloading, resource allocation, and replica
selection. Their central finding is that inter-pillar synergy---jointly optimizing
across all five pillars---remains almost entirely unaddressed in existing work.

The DAG model, formalized in~\cite{zhu2025dag} as $G = (\mathcal{V}, \mathcal{E})$
with per-node properties $\langle \ell_i, d_i^{\text{in}}, d_i^{\text{out}}
\rangle$, enables partial offloading at step granularity: lightweight steps can
execute on fog while GPU-heavy steps run on cloud, with data dependencies
respected. The total application latency~\eqref{eq:dag_latency} now depends on
inter-step data transfer costs, which vary depending on whether consecutive steps
are co-located. This fundamentally changes the optimization problem---instead of
one routing decision per task, there is one decision per DAG step.

\subsection{Deep Reinforcement Learning for Task Offloading}

\subsubsection{DRL for Vehicular Fog Computing}

An anonymous vehicular fog computing study~\cite{vfc_drl2024} trained DQN agents
on a grid-city simulation environment modeled on Salt Lake City's 200$\times$200\,m
block structure. Client vehicles with insufficient compute offload tasks to nearby
parked service vehicles acting as fog nodes. The DQN observes queue depths,
distances, and channel quality at all candidate nodes and selects the best
destination via epsilon-greedy policy. Key result: DRL reduces queue congestion
across the network---not just average per-task latency---and handles high-load
scenarios where greedy methods cascade-fail.

However, this work has no energy objective, no offline pre-training (cold-start
problem), and no mobility handling during task execution---all of which motivate
the proposed architecture.

\subsubsection{SDN-Integrated DRL: CSFD}

The CSFD framework~\cite{csfd2025} embeds a DRL agent directly in the SDN control
plane. Unlike standard SDN controllers that reactively install rules per new flow
(8--15\,ms overhead), CSFD's agent predicts upcoming traffic bursts and
pre-installs OpenFlow rules before the traffic arrives---reducing per-flow overhead
to zero. Evaluated in a VANET environment, CSFD achieves higher packet delivery
ratio and lower latency than both DRL-only and SDN-only baselines. This work
directly motivates the network-routing agent (Agent~2) in the proposed
architecture; however, CSFD handles network routing only, with no compute
placement intelligence---a critical gap.

\subsubsection{DAG Task Offloading with Meta-RL and NSGA-II Pre-training}

Zhu et al.~\cite{zhu2025dag} proposed the architecturally closest published
system to the proposed PCNME. They model applications as DAGs, run NSGA-II
offline on historical batches to generate Pareto-optimal routing plans, convert
these plans into (state, action) training pairs, and pre-train a Meta-RL agent
via behavioral cloning. The meta-policy adapts rapidly to new network environments
via multiple MDPs, eliminating the cold-start problem. Their result: meta-RL
combined with NSGA-II pre-training outperforms both stand-alone DRL and NSGA-II
on dynamic environments.

Critically, this work has no mobility model and no SDN layer; their future work
explicitly requests both---provided by the proposed architecture.

\subsection{Mobility-Aware Offloading}

\subsubsection{Reactive Handoff: NTB/HTB Dual-Buffer System}

A collaborative offloading framework~\cite{paper5_mobility} introduced the
NTB/HTB dual-buffer system as a reactive fallback for mid-computation vehicle
handoffs. When a vehicle disconnects from Fog~A mid-task, the task migrates from
the New Task Buffer (NTB) to the Handoff Task Buffer (HTB), which preempts the
NTB and receives maximum CPU priority. Fog~A completes the task, holds the result
keyed by vehicle~ID, and delivers it the moment the vehicle reconnects to any fog
node---eliminating task re-submission. This achieves 50\%+ latency reduction in
high-mobility scenarios.

However, the system is purely reactive: it waits for physical disconnection before
acting, wasting the entire handoff gap (typically 0.5--3\,s) in wasted service time.

\subsubsection{Proactive Handoff: MAVTO}

MAVTO~\cite{mavto2024} introduced the first peer-reviewed proactive mobility
framework for vehicular edge computing. Given a vehicle's Spatial Tag (GPS
coordinates, speed $v$, heading $\theta$), the system computes:
\begin{equation}
  T_{\text{exit}} = \frac{R_{\text{fog}} - \text{dist}(v_{\text{pos}},\, f_{\text{pos}})}
  {v_{\text{closing}}}
  \label{eq:texit}
\end{equation}
where $v_{\text{closing}}$ is the component of vehicle velocity directed toward
the zone boundary. MAVTO selects among three modes: \emph{direct}
($T_{\text{exec}} < T_{\text{exit}}$---task finishes before vehicle leaves),
\emph{predictive} ($T_{\text{exec}} > T_{\text{exit}}$---pre-spin container at
destination fog node), and \emph{hybrid} (partial local execution). Evaluated on
real vehicular mobility datasets, MAVTO achieves 25\%+ improvement in task
completion success rate over non-mobility-aware baselines.

Limitations: mode selection is rule-based (not learned), no multi-objective
optimization, and no SDN routing for migration paths.

\subsubsection{Spatiotemporal Prediction: Digital Twin + STGNN}

The state of the art in trajectory prediction for vehicular edge computing is
presented by Zhu et al.~\cite{zhu2025spatiotemporal}, who combine a Digital Twin
(continuously synchronized virtual replica of the network), a
Spatio-Temporal Graph Neural Network (STGNN) for demand forecasting, and A3C
multi-agent DRL for task caching. The STGNN models the vehicular network as a
graph $G = (\mathcal{N}, \mathcal{E})$ and applies graph convolution (spatial
correlations) combined with LSTM (temporal patterns) to predict future task
demand at each fog zone, enabling strategic resource pre-allocation 5--10\,s ahead
of actual arrivals.

This 2025 system represents the ceiling of current capabilities but lacks
multi-objective optimization, task pre-filtering, and energy awareness.

\subsection{Comparative Summary}

Table~\ref{tab:comparison} summarizes our ten primary references across six key
dimensions. The proposed architecture (final row) is the only system addressing
all six simultaneously.

\begin{table}[H]
\centering
\caption{Comparative summary of surveyed works.
  \checkmark~=~addressed, $\circ$~=~partial, $\times$~=~not addressed.}
\label{tab:comparison}
\renewcommand{\arraystretch}{1.25}
\begin{tabularx}{\textwidth}{@{}L{3.0cm}C{1.1cm}C{1.2cm}C{1.0cm}C{1.0cm}
                              C{1.0cm}C{1.2cm}C{1.1cm}@{}}
\toprule
\textbf{Work} & \textbf{Multi-obj.} & \textbf{Pre-filter}
  & \textbf{DAG} & \textbf{DRL} & \textbf{SDN}
  & \textbf{Mobility} & \textbf{Year} \\
\midrule
Jafari~\&~Rezvani~\cite{jafari2021}
  & \checkmark & $\times$ & $\times$ & $\times$ & $\times$ & $\times$ & 2021 \\
Nouri et al.~\cite{nouri2024tof}
  & \checkmark & \checkmark & $\times$ & $\times$ & $\times$ & $\times$ & 2024 \\
Gasmi~\&~Ezzedine~\cite{gasmi2024sdn}
  & $\times$ & $\times$ & $\times$ & $\times$ & \checkmark & $\times$ & 2023 \\
Li et al.~\cite{li2025microservice}
  & $\circ$ & $\times$ & \checkmark & $\circ$ & $\times$ & $\times$ & 2025 \\
Collab.\ Framework~\cite{paper5_mobility}
  & $\times$ & $\times$ & $\times$ & $\times$ & $\times$ & $\circ$ & 2023 \\
VFC-DRL~\cite{vfc_drl2024}
  & $\times$ & $\times$ & $\times$ & \checkmark & $\times$ & $\times$ & 2024 \\
CSFD~\cite{csfd2025}
  & $\times$ & $\times$ & $\times$ & \checkmark & \checkmark & $\times$ & 2025 \\
Zhu et al.\ DAG+MRL~\cite{zhu2025dag}
  & \checkmark & $\times$ & \checkmark & \checkmark & $\times$ & $\times$ & 2025 \\
MAVTO~\cite{mavto2024}
  & $\times$ & $\times$ & $\times$ & $\times$ & $\times$ & \checkmark & 2024 \\
Zhu et al.\ STGNN~\cite{zhu2025spatiotemporal}
  & $\times$ & $\times$ & $\circ$ & \checkmark & $\times$ & \checkmark & 2025 \\
\midrule
\textbf{Proposed PCNME}
  & \checkmark & \checkmark & \checkmark & \checkmark
  & \checkmark & \checkmark & 2026 \\
\bottomrule
\end{tabularx}
\end{table}

% =============================================================================
\section{Open Research Gaps}
\label{sec:gaps}
% =============================================================================

The survey reveals five clearly defined gaps. We formally state all five; the
first two constitute our primary contributions to the gap analysis.

\medskip\noindent\textbf{Gap~1 — The Integration Gap (Primary).}
No published system simultaneously addresses all six dimensions in
Table~\ref{tab:comparison}: multi-objective optimization, pre-filtering,
DAG task modeling, DRL online adaptation, SDN network management, and proactive
mobility handling. Systems that advance one or two dimensions consistently
ignore the others. The lack of integration is not incidental---it reflects the
absence of a unified theoretical framework that models the interactions between
compute placement, network routing, and mobility as a single joint optimization
problem. Specifically, compute placement decisions (which fog node runs a DAG
step) and network routing decisions (which physical path carries the step's data)
are interdependent: a placement that minimizes compute latency may select a fog
node whose network path is currently saturated, negating the benefit. Resolving
this gap requires either joint optimization (intractable action space) or
architectural coupling of two specialized agents with shared information---a
design not yet demonstrated at the DAG-and-mobility scale.


\medskip\noindent\textbf{Gap~2 — The Cold-Start Gap (Primary).}
All reviewed DRL systems are trained from random weight initialization, producing
near-random policies during early deployment. This cold-start period---spanning
thousands to tens of thousands of interaction steps---causes significant service
degradation that is unacceptable in production vehicular safety systems. Zhu
et al.~\cite{zhu2025dag} demonstrate NSGA-II Pareto solutions as behavioral
cloning data for meta-RL pre-training, but do not validate this pipeline in the
full IoT--Fog--Cloud vehicular context with mobility and SDN. The critical
open question is: can offline evolutionary optimization (TOF-MMDE-NSGA-II)
systematically eliminate the cold-start problem for both compute placement and
network routing agents, and does this pre-training advantage persist across the
diverse environmental conditions of a city-scale vehicular network?


\medskip\noindent\textbf{Gap~3 — The DAG-Mobility Gap.}
Mobility-aware offloading systems~\cite{mavto2024,paper5_mobility} treat tasks
as monolithic units. When a task is structured as a multi-step DAG, a handoff
may occur between any two successive steps, requiring step-level migration
decisions. No published work addresses proactive handoff for DAG-structured
workloads where different steps may execute concurrently on different fog nodes.


\medskip\noindent\textbf{Gap~4 — The Pebble Avalanche Gap.}
The TOF-Broker~\cite{nouri2024tof} resolves the boulder problem but introduces
the Pebble Avalanche: under high task arrival rates, the pebble queue grows
faster than NSGA-II drains it, causing cascade failures. No mechanism for
graceful degradation under avalanche conditions (e.g., dynamic threshold
adjustment or task aggregation and bulk cloud routing) has been proposed or
validated.


\medskip\noindent\textbf{Gap~5 — The Spatiotemporal Prediction Gap.}
MAVTO~\cite{mavto2024} uses linear dead-reckoning for trajectory prediction,
which fails under sudden maneuvers, intersections, and variable speeds. Zhu
et al.~\cite{zhu2025spatiotemporal} demonstrate STGNN-based spatiotemporal
demand prediction as a superior alternative, but do not integrate it with a
proactive handoff mechanism or multi-objective optimization. A system combining
STGNN prediction with proactive container migration and Pareto-optimal placement
has not been demonstrated.


%% TODO-TEAMMATE-1: Insert your 2+ research gaps here, derived from your
%% own 10 papers. Use the ...
 environment to
%% match the formatting. Label them Gap~6, Gap~7, etc.
%%
%% Example structure:
%% %% \medskip\noindent\textbf{Gap~6 — [Your Gap Title].}
%% [Your gap description citing your papers with \cite{your_key}.]
%%


%% TODO-TEAMMATE-2: Same as above. Insert Gap~7, Gap~8 etc.

% =============================================================================
\section{Proposed Architecture: Predictive Cloud-Native Mobile Edge}
\label{sec:proposed}
% =============================================================================

\subsection{Architecture Overview}

The Predictive Cloud-Native Mobile Edge (PCNME) architecture directly addresses
Gaps~1--5 through five integrated components:

\begin{enumerate}[leftmargin=*]
  \item \textbf{TOF-Broker + Aggregator (Gaps~1, 4).}
    The TOF-Broker applies EC classification~\eqref{eq:ec} per DAG step, routing
    boulders directly to cloud. The Aggregator handles avalanche conditions by
    compressing overflow pebbles into Super-Tasks routed via SDN VIP lanes.

  \item \textbf{TOF-MMDE-NSGA-II Offline Optimizer (Gap~2).}
    Combines TOF pre-filtering with MMDE-enhanced NSGA-II~\eqref{eq:mmde} to
    generate Pareto-optimal routing plans from historical data. Plans are
    decomposed into (state, action) training pairs for behavioral cloning of
    both DRL agents, eliminating cold-start.

  \item \textbf{DRL Agent~1 --- Task Placement Brain (Gap~1).}
    A Deep Q-Network~\cite{vfc_drl2024} pre-trained on NSGA-II Pareto solutions
    and fine-tuned online via experience replay. Observes a 13-dimensional state
    vector (fog CPU loads, queue depths, task EC, bandwidth, vehicle speed,
    heading, $T_{\text{exit}}$, deadline remaining, cloud load). Action space:
    \{FOG\_A, FOG\_B, FOG\_C, FOG\_D, CLOUD\}.

  \item \textbf{DRL Agent~2 --- SDN Network Brain (Gaps~1, 4).}
    A second DQN embedded in the SDN controller~\cite{csfd2025}, pre-trained on
    historical city traffic patterns. Pre-installs OpenFlow flow rules before
    traffic arrives, eliminating per-flow controller overhead. Decoupled from
    Agent~1 but shares network-state telemetry.

  \item \textbf{Trajectory Predictor + NTB/HTB System (Gaps~3, 5).}
    Parses vehicle Spatial Tags (GPS, speed, heading) and computes
    $T_{\text{exit}}$~\eqref{eq:texit} per DAG step individually---step-level
    mobility assessment, not task-level. Triggers proactive container pre-spin
    on the destination fog node when $T_{\text{exec}} > T_{\text{exit}}$.
    NTB/HTB buffers~\cite{paper5_mobility} provide deterministic fallback when
    trajectory prediction fails.
\end{enumerate}

\subsection{Hybrid Offline--Online AI Pipeline}

\begin{enumerate}[leftmargin=*, label=\textbf{Phase \arabic*:}]
  \item \textbf{Offline.} Historical task logs and mobility traces feed
    TOF-MMDE-NSGA-II. Pareto solutions are converted to (state, action) pairs
    and used to pre-train both agents via behavioral cloning.
  \item \textbf{Online.} Agent~1 makes placement decisions per pebble DAG step
    at inference speed ($<$1\,ms). Agent~2 manages SDN flow rules proactively.
    Both agents update weights continuously from simulation rewards:
    \[
      R = -0.5\cdot L_{\text{norm}} - 0.3\cdot E_{\text{norm}}
        - 0.2\cdot \mathbb{1}[L > D] \cdot \lambda_{\text{penalty}}
    \]
    where $\lambda_{\text{penalty}} = 10$ for safety-critical DAG steps.
\end{enumerate}

% =============================================================================
\section{Comparative Discussion}
\label{sec:discussion}
% =============================================================================

\subsection{Against Generation-2 Systems}

The proposed architecture subsumes~\cite{jafari2021} entirely: NSGA-II with MMDE
is retained as the offline optimizer. The TOF-Broker (from~\cite{nouri2024tof})
reduces the problem NSGA-II must solve; the DRL agents provide real-time
adaptation that NSGA-II cannot; the mobility layer handles vehicular dynamics
neither predecessor addressed.

\subsection{Against DAG+MRL \cite{zhu2025dag}}

Zhu et al.'s meta-RL system is the architecturally closest prior work. The
proposed system extends it with: (i)~TOF pre-filtering reducing Agent~1's
decision space, (ii)~SDN-managed network routing via Agent~2, and
(iii)~step-level proactive mobility via~$T_{\text{exit}}$~\eqref{eq:texit}.
Their future work explicitly requests all three.

\subsection{Against CSFD \cite{csfd2025}}

CSFD provides state-of-the-art SDN+DRL network routing but has no compute
placement agent. The proposed system's Agent~2 performs the same function as
CSFD; Agent~1 provides the compute placement intelligence CSFD lacks. The
architectural coupling between agents addresses Gap~1 in a way CSFD cannot.

\subsection{Against MAVTO \cite{mavto2024}}

MAVTO provides the $T_{\text{exit}}$ proactive handoff mechanism but applies it
to monolithic tasks. The proposed system extends MAVTO with: (i)~step-level
assessment (Gap~3), (ii)~RL-learned mode selection replacing MAVTO's rule-based
thresholds, and (iii)~SDN-managed migration paths for container transfer.

%% TODO-TEAMMATE-1: Add a subsection here comparing your proposed ideas
%% (from your papers) against the state of the art in your sub-domain.
%% Example: \subsection{Against [Your Area] Systems} ...

%% TODO-TEAMMATE-2: Same — add your comparative subsection here.

% =============================================================================
%% TODO-TEAMMATE-1: INSERT YOUR FULL SECTION HERE
%% ─────────────────────────────────────────────────────────────────────────────
%% This is where your teammate adds their literature review section.
%% Suggested structure:
%%
%% \section{[Your Research Sub-Domain]}
%% \label{sec:teammate1}
%%
%% \subsection{Background}
%% ...
%%
%% \subsection{State of the Art}
%% ...
%%
%% \subsection{Open Problems}
%% ...
%%
%% Cite your papers using \cite{} and add them to the bibliography at the end.
%% ─────────────────────────────────────────────────────────────────────────────

% =============================================================================
%% TODO-TEAMMATE-2: INSERT YOUR FULL SECTION HERE
%% ─────────────────────────────────────────────────────────────────────────────
%% Same as above for the second teammate.
%% ─────────────────────────────────────────────────────────────────────────────

% =============================================================================
\section{Conclusion}
\label{sec:conclusion}
% =============================================================================

This literature review traced the evolution of task offloading research across
four generations---from single-objective heuristics, through multi-objective
evolutionary algorithms and hybrid pre-filtering systems, to the current generation
of adaptive AI-driven architectures. We surveyed ten representative works spanning
2021--2025 and formally identified five open research gaps.

Gaps~1 and~2---the Integration Gap and Cold-Start Gap---represent our primary
contributions to the gap analysis. Gap~1 establishes that the joint optimization
of compute placement, network routing, and mobility under DAG task models has
not been demonstrated at city scale. Gap~2 argues that the use of offline
evolutionary Pareto solutions as pre-training data for online DRL agents is a
promising and largely unvalidated approach to eliminating cold-start degradation
in vehicular safety systems.

The proposed PCNME architecture addresses all five gaps through a unified pipeline:
TOF-MMDE-NSGA-II offline optimization $\rightarrow$ dual DQN agent pre-training
$\rightarrow$ live task placement and SDN routing with proactive GPS-driven
mobility handling. No individual component is novel in isolation; the novelty
lies in their integration and the offline-to-online knowledge transfer mechanism.

Future work directions include: (1) validating the behavioral cloning pre-training
pipeline across diverse city topologies; (2) extending trajectory prediction from
linear dead-reckoning to STGNN-based spatiotemporal
forecasting~\cite{zhu2025spatiotemporal}; (3) investigating federated learning
to train Agent~2 across multiple city deployments without sharing raw traffic
data; and (4) exploring TD3 or SAC as higher-performing alternatives to DQN for
the continuous-action resource allocation sub-problem.

%% TODO-TEAMMATE-1: Add 1-2 sentences summarizing your conclusion here.
%% TODO-TEAMMATE-2: Add 1-2 sentences summarizing your conclusion here.

% =============================================================================
% BIBLIOGRAPHY
% =============================================================================
\bibliographystyle{plainnat}

\begin{thebibliography}{99}

% ── YOUR 10 PAPERS ──────────────────────────────────────────────────────────

\bibitem{jafari2021}
Jafari, V. and Rezvani, M.\,H.
\newblock Joint optimization of energy consumption and time delay in
  {IoT}--fog--cloud computing environments using {NSGA-II} metaheuristic
  algorithm.
\newblock \emph{Journal of Ambient Intelligence and Humanized Computing},
  14:1675--1698, 2023 (published online 2021).

\newblock URL: \url{https://link.springer.com/article/10.1007/s12652-021-03388-2}

\bibitem{nouri2024tof}
Nouri, N. and Rezvani, M.\,H.
\newblock A multi-objective approach for optimizing {IoT} applications
  offloading in fog--cloud environments with {NSGA-II}.
\newblock \emph{The Journal of Supercomputing}, 2024.

\newblock URL: \url{https://link.springer.com/article/10.1007/s11227-024-06431-z}

\bibitem{gasmi2024sdn}
Gasmi, K. and Ezzedine, T.
\newblock Reinforcement learning based task offloading of {IoT} applications in
  fog computing: algorithms and optimization techniques.
\newblock \emph{Cluster Computing}, 2024.

\newblock URL: \url{https://www.ijcnis.org/index.php/ijcnis/article/view/8458/2522}

\bibitem{li2025microservice}
Li, Y. et al.
\newblock Energy-efficient resource management in microservices-based fog and
  edge computing: state-of-the-art and future directions.
\newblock \emph{ACM Computing Surveys}, 2025.
\newblock URL: \url{https://arxiv.org/pdf/2512.04093}

\bibitem{paper5_mobility}
{Anonymous}.
\newblock A collaborative offloading task framework for {IoT} fog computing.
\newblock Technical report, 2023.

\bibitem{vfc_drl2024}
{Anonymous}.
\newblock Deep reinforcement learning for delay-optimized task offloading in
  vehicular fog computing.
\newblock \emph{arXiv preprint arXiv:2410.03472}, 2024.
\newblock URL: \url{https://arxiv.org/html/2410.03472}

\bibitem{csfd2025}
{Anonymous}.
\newblock {CSFD}: Enhancing congestion reduction in vehicular communication
  with {SDN} and fog integration using deep reinforcement learning.
\newblock \emph{Analog Integrated Circuits and Signal Processing}, 2025.

\newblock URL: \url{https://link.springer.com/article/10.1007/s10470-025-02541-7}

\bibitem{zhu2025dag}
Zhu, L. et al.
\newblock Dynamic task offloading scheme for edge computing via
  meta-reinforcement learning.
\newblock \emph{Journal of Manufacturing Systems}, 2025.
\newblock URL: \url{https://www.sciencedirect.com/science/article/pii/S1546221825000712}

\bibitem{mavto2024}
{Anonymous}.
\newblock Mobility-aware task offloading and resource allocation in
  {UAV}-assisted vehicular edge computing networks.
\newblock \emph{Drones}, 8(11):696, 2024.

\newblock URL: \url{https://www.mdpi.com/2504-446X/8/11/696}

\bibitem{zhu2025spatiotemporal}
Zhu, L. et al.
\newblock A vehicular edge computing offloading and task caching solution based
  on spatiotemporal prediction.
\newblock \emph{Future Generation Computer Systems}, 166:107679, 2025.

\newblock URL: \url{https://www.sciencedirect.com/article/abs/pii/S0167739X24006435}

% ── ADDITIONAL REFERENCES USED IN TEXT ──────────────────────────────────────

\bibitem{mashal2024drl}
Mashal, I. et al.
\newblock Multiobjective offloading optimization in fog computing using deep
  reinforcement learning.
\newblock \emph{Journal of Computer Networks and Communications}, 2024.


%% TODO-TEAMMATE-1: Add your 10 paper references here using the same
%% \bibitem format. Give each a short memorable key like {author2024keyword}.

%% TODO-TEAMMATE-2: Add your 10 paper references here.

\end{thebibliography}

\end{document}
